\begin{proposition}[DP29: a valid scalar estimate and a sharp nonzero example]\label{DP29}
Under (H2)--(H4), let $U\in R[[S_j:j\in J]]$, with $J$ finite, let $\ell_j(\mathbf t)$ be integral linear forms in $N$ variables, let $\epsilon\in R^\times$, and let $\sigma\in\{1,-1\}$. Set
\[
 \Theta(\mathbf t)=U(E_h(\ell_j(\mathbf t)):j\in J)
 \exp\left(\frac\sigma2\sum_j\beta(\ell_j(\mathbf t),\epsilon\ell_j(\mathbf t))\right).
\]
For $n=|\kappa|$,
\begin{equation}\label{eq:scalar-old-residual}
 v\bigl(\kappa![\mathbf t^\kappa]\Theta\bigr)
 \ge \rho_2n-\frac{\sum_i s_p(\kappa_i)}{p-1}
 \ge\rho_2n-N(\log_p n+1)\quad(n\ge1).
\end{equation}
For the nonzero correction $\beta$ in DP28, the particular function
\[
 \Theta_h(t)=E_h(t)\left(1+E_h(t)^pE_h(\epsilon t)\right)^{1/2}
\]
has, for every $n=q^r$,
\begin{equation}\label{eq:sharp-corrected-scalar}
 v(n![t^n]\Theta_h)=v_p(n!)-\frac{n-1}{q-1}=c_h(n-1).
\end{equation}
Thus the failed intermediate estimate is compatible with exact nonvanishing at the larger derivative slope.
\end{proposition}
\begin{proof}
Write $w_r(A)=\inf_\nu(v(A_\nu)+r|\nu|)$. Products add lower weight bounds, and substitution of series of weight at least $r$ preserves them; these assertions hold coefficient by coefficient. Formula \eqref{eq:honda-scale} gives $w_r(E_h)\ge r$ for $r\ge\rho_h$.

Torus invariance and positive bidegrees force $\ord\beta\ge p+1$: if $a,b>0$ and $a+b\le p$, then $a-pb\ne0$ and $|a-pb|<p^2-1\le q-1$. Put $r_0=p/(p^2-1)$. Then $r_0\ge\rho_h$ and $(p+1)r_0=s+1$. Both integral series $H_X-1,H_Y-1$ have order at least $p+1$. Substituting $E_h$ gives weight at least $s+1$, so their formal logarithms have the same lower bound: the $m$-th term has weight at least $m(s+1)-v_p(m)\ge s+1$. Hence
\[
 w_{r_0}(\delta_X\beta),\ w_{r_0}(\delta_Y\beta)\ge s+1.
\]
For every integer $a\ge2$ there is $0<i<a$ with $v_p\binom ai\le1$. If $a$ is not a power of $p$, expanding $(1+T)^a=\prod_r(1+T^{p^r})^{a_r}$ modulo $p$ supplies an interior $i=p^r$ with unit coefficient. If $a=p^r$, take $i=p^{r-1}$ and use $\binom{p^r}{p^{r-1}}=p\binom{p^r-1}{p^{r-1}-1}$, the second factor being a unit. Every monomial of $\beta$ has one exponent at least two. Its corresponding coefficient in one coboundary is exactly this binomial coefficient times $\beta_{a,b}$. Thus $w_{r_0}(\beta)\ge s$.

Integral linear substitutions and the factor $\sigma/2$ preserve this bound, since $p$ is odd. The exponential has weight at least zero by $v_p(m!)\le ms$; the $U$ factor does also. Therefore $v([\mathbf t^\kappa]\Theta)\ge-r_0n$. Adding Legendre's exact formula for $v_p(\kappa!)$ gives \eqref{eq:scalar-old-residual}, since $s-r_0=\rho_2$. The digit estimate gives its final inequality; zero coordinates contribute zero. At $n=0$ only integrality of $U(0)$ is required.

For the explicit example, $\binom{1/2}{m}\in\Z_p$ for odd $p$: approximate $1/2$ by nonnegative integers in $\Z_p$ and use continuity of each binomial polynomial. Hence
\[
 \pi^{-1}\Theta_h(\pi z)
 =E_{h,\pi}(z)\left(1+\pi^{p+1}E_{h,\pi}(z)^pE_{h,\pi}(\epsilon z)\right)^{1/2}
 \equiv E_{h,\pi}(z)\pmod\pi.
\]
Its coefficient at $z^{q^r}$ is a unit by \eqref{eq:honda-scale}. Rescaling and $s_p(q^r)=1$ prove \eqref{eq:sharp-corrected-scalar}.
\end{proof}
